% Figure 1: study structure. Plain block layout. No color, no dashed lines, few arrows.
\begin{figure*}[t]
\centering
\resizebox{\textwidth}{!}{%
\begin{tikzpicture}[
  font=\sffamily,
  cell/.style={rectangle, draw=black, minimum height=1.15cm, minimum width=4.0cm,
               align=center, inner sep=4pt},
  band/.style={rectangle, draw=black, thick, minimum height=1.15cm, minimum width=16.6cm,
               align=center, inner sep=4pt},
  arr/.style={-{Stealth[length=6pt]}, thick},
]

% Row 1: the four axes that we vary
\node[cell] (a1) at (0,0)     {\textbf{Models}\\{\small 8 commercial, 4 open}};
\node[cell] (a2) at (4.2,0)   {\textbf{Hardware}\\{\small RTX PRO 6000, H100, A6000}};
\node[cell] (a3) at (8.4,0)   {\textbf{Benchmarks}\\{\small 4 suites, 424 tasks}};
\node[cell] (a4) at (12.6,0)  {\textbf{Budget}\\{\small samples, rounds, turns}};

% Row 2: the harness that stays fixed
\node[band] (harness) at (6.3,-2.15)
  {\textbf{One fixed evaluation harness}\\[2pt]
   {\small build prompt \ \textbar\ \ generate \ \textbar\ \ parse code \ \textbar\ \
    \textbf{score in an isolated subprocess} \ \textbar\ \ classify error \ \textbar\ \ repair and retry}};

% Row 3: what we read out
\node[band] (out) at (6.3,-4.15)
  {\textbf{What we measure}\\[2pt]
   {\small correctness, Fast-0 and Pass@1 \ \textbar\ \ speedup, Fast-1 and Fast-2 \ \textbar\ \
    error types \ \textbar\ \ wall-clock and completion rate}};

% Row 4: research questions
\node[band] (rq) at (6.3,-6.15)
  {\textbf{Research questions}\\[2pt]
   {\small RQ1 how accurate \ \textbar\ \ RQ2 how much the setup decides \ \textbar\ \
    RQ3 how failures look \ \textbar\ \ RQ4 what raises accuracy \ \textbar\ \ RQ5 what transfers}};

% Three arrows only
\draw[arr] (6.3,-0.62) -- (harness.north);
\draw[arr] (harness.south) -- (out.north);
\draw[arr] (out.south) -- (rq.north);

\end{tikzpicture}%
}
\caption{Structure of the study. We keep one evaluation harness fixed and change four axes on top of it.
The harness scores every kernel in an isolated subprocess. This is required, not optional, because one
bad kernel can corrupt the CUDA context, and then every later evaluation in the same process fails.
We do not propose a new generation system. We hold the harness constant and change one axis at a time,
so that we can tell apart the effect of the model, the hardware, the benchmark, and the compute budget.}
\label{fig:overview}
\end{figure*}
